% =========================================================
% Oscillation — Canonical Notes
% Chapter: continuity
% Volume III · Analysis
% =========================================================

\section{Oscillation}
\label{sec:oscillation}

% ---------------------------------------------------------
% Toolkit
% ---------------------------------------------------------

\begin{tcolorbox}[title={Toolkit — Oscillation}]
\begin{itemize}
  \item $\omega(f; E) = \sup_{x_1, x_2 \in E} |f(x_1) - f(x_2)| = \operatorname{diam}(f(E))$.
  \item $\omega(f; x_0) = \lim_{\delta \to 0^+} \omega(f; U_\delta(x_0) \cap E)$.
  \item The limit defining $\omega(f; x_0)$ exists because the neighborhood oscillation is monotone nonincreasing in $\delta$.
  \item $f$ continuous at $x_0$ $\iff$ $\omega(f; x_0) = 0$.
  \item Discontinuity set: $D(f) = \{x \in E : \omega(f; x) > 0\} = \bigcup_{n=1}^{\infty}\{x \in E : \omega(f; x) \geq 1/n\}$.
\end{itemize}
\end{tcolorbox}

% ---------------------------------------------------------
% Definition — Oscillation on a Set
% ---------------------------------------------------------

\begin{tcolorbox}[title={Definition — Oscillation on a Set}]
\begin{definition}[Oscillation on a Set]
\label{def:oscillation-on-a-set}
Let $E \subseteq \mathbb{R}$ and let $f : E \to \mathbb{R}$. The
\textbf{oscillation of $f$ on $E$} is
\[
  \omega(f; E)
  \;:=\;
  \sup_{x_1, x_2 \in E} |f(x_1) - f(x_2)|
  \;=\;
  \operatorname{diam}(f(E)).
\]
\end{definition}
\end{tcolorbox}

\begin{remark*}[Standard quantified statement]
\[
  \omega(f; E)
  \;=\;
  \sup\bigl\{\, |f(x_1) - f(x_2)| : x_1, x_2 \in E \,\bigr\}.
\]
\end{remark*}

\begin{remark*}[Interpretation]
The oscillation on $E$ measures the diameter of the image $f(E)$: the
largest separation between any two function values on the set. It does
not record where the extreme values occur, only how widely the image
spreads. Oscillation zero means $f$ is constant on $E$.
\end{remark*}

% ---------------------------------------------------------
% Definition — Oscillation at a Point
% ---------------------------------------------------------

\begin{tcolorbox}[title={Definition — Oscillation at a Point}]
\begin{definition}[Oscillation at a Point]
\label{def:oscillation-at-a-point}
Let $E \subseteq \mathbb{R}$, let $f : E \to \mathbb{R}$, and let
$x_0 \in E$. The \textbf{oscillation of $f$ at $x_0$} is
\[
  \omega(f; x_0)
  \;:=\;
  \lim_{\delta \to 0^+} \omega\bigl(f;\, U_\delta(x_0) \cap E\bigr),
\]
where $U_\delta(x_0) = \{x \in \mathbb{R} : |x - x_0| < \delta\}$.
\end{definition}
\end{tcolorbox}

\begin{remark*}[Standard quantified statement]
\[
  \omega(f; x_0)
  \;=\;
  \lim_{\delta \to 0^+}
  \sup\bigl\{\, |f(x_1) - f(x_2)| : x_1, x_2 \in U_\delta(x_0) \cap E \,\bigr\}.
\]
\end{remark*}

\begin{remark*}[The limit exists]
For $0 < \delta_1 < \delta_2$, the inclusion
$U_{\delta_1}(x_0) \cap E \subseteq U_{\delta_2}(x_0) \cap E$
forces
\[
  \omega\bigl(f;\, U_{\delta_1}(x_0) \cap E\bigr)
  \;\leq\;
  \omega\bigl(f;\, U_{\delta_2}(x_0) \cap E\bigr).
\]
The function $\delta \mapsto \omega(f; U_\delta(x_0) \cap E)$ is
therefore monotone nonincreasing as $\delta \to 0^+$ and bounded
below by $0$. The limit exists.
\end{remark*}

\begin{remark*}[Interpretation]
The pointwise oscillation measures the residual spread of nearby
function values as the neighborhood shrinks to zero. If the nearby
values of $f$ collapse to a single cluster, $\omega(f; x_0) = 0$.
If a positive amount of spread persists no matter how tightly the
neighborhood is drawn, $\omega(f; x_0) > 0$.
\end{remark*}

% ---------------------------------------------------------
% Theorem — Continuity Characterised by Oscillation
% ---------------------------------------------------------

\begin{theorem}[Continuity Characterised by Oscillation]
\label{thm:continuity-iff-zero-oscillation}
Let $E \subseteq \mathbb{R}$, let $f : E \to \mathbb{R}$, and let
$x_0 \in E$. Then
\[
  f \text{ is continuous at } x_0
  \;\iff\;
  \omega(f; x_0) = 0.
\]
\end{theorem}

\begin{remark*}[Standard quantified statement]
\[
\begin{aligned}
&\Bigl(\forall \varepsilon > 0,\;
\exists \delta > 0,\;
\forall x \in E :\;
|x - x_0| < \delta \Rightarrow |f(x) - f(x_0)| < \varepsilon\Bigr) \\
&\qquad\iff\qquad
\lim_{\delta \to 0^+}
\omega\bigl(f;\, U_\delta(x_0) \cap E\bigr) = 0.
\end{aligned}
\]
\end{remark*}

\begin{remark*}[Negated quantified statement]
\[
  f \text{ discontinuous at } x_0
  \;\iff\;
  \omega(f; x_0) > 0.
\]
\end{remark*}

\begin{remark*}[Interpretation]
Continuity at $x_0$ is equivalent to the local spread of $f$
collapsing to zero. The forward direction follows because, if $f$
is continuous at $x_0$, every nearby value $f(x)$ is close to
$f(x_0)$, hence any two nearby values are close to each other.
The backward direction reverses this: if $\omega(f; x_0) = 0$, then
on sufficiently small neighborhoods all function values are confined
to an arbitrarily short interval containing $f(x_0)$.

This criterion replaces the $\varepsilon$-$\delta$ alternation with
a single local invariant. Continuity becomes a local rigidity
statement: no oscillation survives arbitrarily close to the point.
\end{remark*>

% ---------------------------------------------------------
% Proposition — Discontinuity Set via Oscillation
% ---------------------------------------------------------

\begin{proposition}[Discontinuity Set via Oscillation]
\label{prop:discontinuity-set-via-oscillation}
Let $E \subseteq \mathbb{R}$ and let $f : E \to \mathbb{R}$. The
discontinuity set of $f$ satisfies
\[
  D(f)
  \;=\;
  \{x \in E : \omega(f; x) > 0\}
  \;=\;
  \bigcup_{n=1}^{\infty}
  \left\{x \in E : \omega(f; x) \geq \tfrac{1}{n}\right\}.
\]
\end{proposition}

\begin{remark*}[Standard quantified statement]
\[
  \forall x \in E :\;
  x \in D(f)
  \;\iff\;
  \omega(f; x) > 0.
\]
\end{remark*}

\begin{remark*}[Interpretation]
The continuity characterisation by oscillation identifies
$D(f)$ with the set where $\omega(f; \cdot)$ is strictly positive.
The Archimedean step $r > 0 \iff \exists n \in \mathbb{N},\; r \geq 1/n$
then yields the union decomposition.

Each threshold set $\{x \in E : \omega(f; x) \geq 1/n\}$ is
closed in $E$. Consequently $D(f)$ is an $F_\sigma$ set — a
countable union of closed sets — which governs the topological
structure of discontinuities and the Riemann integrability criterion.
For bounded functions, Riemann integrability is equivalent to
$D(f)$ having Lebesgue measure zero. Oscillation is the quantity
that makes this criterion transparent.
\end{remark*>
